\section{CDSM}
\label{sec:method}

CDSM's primary policy F runs up to four local \emph{scouts}, each capped
at 400 new distance evaluations, then continues from their reintegrated
state. Local thresholds let routes develop before competing with the
main answers. We describe squared L2, where smaller distances are better.

\begin{figure*}[t]
\centering
\includegraphics[width=\textwidth]{method-feedback.pdf}
\Description{Four stages show a paused main search, bounded local scouts, reintegration of active, deferred and partly expanded nodes, and main continuation. All stages share scores and traversal progress; the answer collector is distinct from executable frontier state.}
\caption{CDSM-F's execution organization (conceptual, not a measured trajectory). Scouts return executable bottom-layer work to the main frontier. The suspended main current node resumes first; the merged candidates then compete by query distance, potentially redirecting later expansions. All new distance evaluations share the query limit $B_q$.}
\label{fig:method}
\end{figure*}

\paragraph{Shared state.}
One workspace shares a distance cache, global top-$k$ collector $A$,
bottom-layer discovery/completion marks, and each opened node's next-neighbor
cursor. Discovery means entering bottom traversal, distinct from being
scored during upper descent. Main frontier $M$ retains active and deferred
candidates plus an optional partially expanded current node. Each scout
has a separate empty local frontier and top-$k$ set, while sharing the
workspace. Every first physical score updates $A$: an upper-layer
discovery can improve the answer without subsequent bottom exploration.

\paragraph{Initial search and entrances.}
The main search descends from the graph's ordinary HNSW entry point.
It spends at most $P_{\max}=6400$ evaluations, including descent.
Between node expansions, it can stop when the closest pending candidate's
squared distance is at least $(1+\gamma)^2 \tau_A$, where $\gamma=0.075$
and $\tau_A$ is the current global $k$th distance.
This condition applies only after $A$ contains $k$ answers.
The search retains its queues and any incomplete adjacency cursor.
Let $p$ denote its actual evaluation count.

A graph-specific table supplies up to 64 distinct upper-layer nodes in a
query-independent, sampled order; imported graphs retain their saved table.
Appendix~\ref{app:execution} gives the exact construction.
The order is not a learned utility ranking.
At each scout, F takes the first table node not previously used as a scout
entrance and not already discovered by bottom-layer traversal.

\paragraph{Local exploration.}
For current distance count $d_c$, a scout receives the boundary
$L=\min(B_q,d_c+400)$.
Its upper descent begins at the selected node's highest layer and charges
all new distances against $L$.
After descent, the local top-$k$ set $T$ is cleared.
If budget remains and the endpoint is not bottom-layer discovered, the
endpoint seeds the local frontier and $T$.
Thus upper-layer scores can improve $A$ without prematurely tightening
the scout's local threshold.

The scout resumes a partially expanded node before selecting another one.
For an undiscovered neighbor, it reads the pre-update threshold
$\tau=T_k$, or $\infty$ while $T$ is not full, and obtains its cached
or newly computed distance $d$.
It marks the neighbor discovered.
If $d<\tau$, it inserts the neighbor into its active queue and updates $T$;
otherwise it retains the neighbor in its deferred queue.
Between expansions, deferred candidates with distance below the current
$\tau$ become active, and completed nodes are discarded.
The scout pauses when no eligible active candidate remains, when the
best active distance exceeds $\tau$, or when $L$ is reached.
A candidate farther than the global $k$th answer can still be locally
promising and expose neighbors absent from the main route.

\paragraph{Reintegration and continuation.}
F inserts each scout's active, deferred, and unfinished current work into
the main active queue, retaining shared cursors, then clears the local
frontier for reuse. Scouts are not independently scheduled afterward.
Continuation resumes the main suspended current node first, then orders
merged candidates by distance. It skips completed nodes, admits new
bottom discoveries without a local threshold, and disables initial
distance stopping. Thus locally rejected work remains available for
further discovery; reintegration retains more than scout answers.

\paragraph{Budgeted execution.}
Algorithm~\ref{alg:cdsm} summarizes the stages.
The total limit is either a fixed $B_q$, known before initialization,
or $B_q=p+E$ under an additional-allowance contract.
With a fixed limit, the initial stage is also capped by $B_q$.

\begin{algorithm}[t]
\small
\begin{algorithmic}[1]
\Require Query $q$, graph $G$, entrance table $\mathcal E$, budget contract
\State Initialize shared workspace $W$, answers $A$, frontier $M$
\State $b_0\gets\min(6400,B)$ if absolute; otherwise $6400$
\State Descend and run main search under $b_0$, $\gamma=0.075$
\State Preserve $M$; $p\gets d_c$ \Comment{actual distance count}
\State $B_q\gets B$ if absolute; otherwise $p+E$
\For{$i=1,\ldots,4$ while $d_c<B_q$}
  \State $e\gets$ first unused, bottom-undiscovered node in $\mathcal E$
  \If{no such $e$} \State \textbf{break} \EndIf
  \State Mark $e$ used; $L\gets\min(B_q,d_c+400)$
  \State $v\gets\Call{UpperDescend}{e,L,W}$
  \State Clear local frontier $S$ and local top-$k$ set $T$
  \If{$d_c<L$ and $v$ is bottom-undiscovered}
    \State Seed $v$ into $S,T$; mark $v$ bottom-discovered
    \State $\Call{LocalExplore}{S,T,L,W}$ using the local rules
  \EndIf
  \State Insert $S$.active, $S$.deferred, $S$.current into $M$.active
  \State Retain shared cursors and $M$.current; clear $S$
\EndFor
\State Resume $M$.current, then merged frontier under $B_q$
\State Use unrestricted admission and no distance stopping
\State \Return $A$
\end{algorithmic}
\caption{CDSM-F. Upper descent and local exploration check their limit before each new score. Final continuation stops at the limit or when recoverable work is exhausted.}
\label{alg:cdsm}
\end{algorithm}

Traversal routines check the relevant boundary before every new scoring
operation. Cache hits do not spend a physical evaluation.
If scout $i$ spends $s_i$ and final continuation spends $c$, then
$p+\sum_i s_i+c\leq B_q$.
Early scout stops leave more budget for continuation; they do not enlarge
another scout's 400-evaluation limit.
F neither adds filler work nor guarantees that every graph permits
exhausting $B_q$.
Equal actual evaluations therefore require measurement, and do not imply
equal queue work or latency.

Paid-selection and scheduling extensions are evaluated separately in
Section~\ref{sec:family}; their selection work uses the same query budget.

\paragraph{Reference configuration and parameter meaning.}
Table~\ref{tab:parameters} fixes the retained development settings for T2I
and transfer, without claiming optimality or a newly preregistered selection.
We test other graphs, allowances, and neighboring scout depths; dimension-
or size-based parameter scaling remains unvalidated.

\begin{table}[t]
\centering\small
\caption{Frozen F parameters and query-dependent quantities.}
\label{tab:parameters}
\begin{tabularx}{\columnwidth}{@{}llY@{}}
\toprule
Symbol & Value & Meaning \\
\midrule
$k$ & 10 & Returned answers; local threshold size \\
$P_{\max}$ & 6,400 & Initial upper and bottom scoring limit \\
$\gamma$ & 0.075 & Initial stopping slack; disabled in final phase \\
$m,s$ & 4, 400 & Maximum scouts and new distances per scout \\
$|\mathcal E|$ & $\leq64$ & Distinct upper-layer entrance table \\
$p$ & Measured & Actual initial expenditure for this query \\
$B_q$ & $B$ or $p+E$ & Total allowed new distances \\
$\tau_A,\tau$ & Dynamic & Global and scout-local $k$th distances \\
\bottomrule
\end{tabularx}
\end{table}

\begin{table*}[t]
\centering\small
\caption{Search organization. C and MEP are controlled policies; FANNG is our same-graph search reconstruction. iQAN is a literature comparison, not a measured baseline in this study.}
\label{tab:semantics}
\begin{tabularx}{\textwidth}{@{}lYYY@{}}
\toprule
Method & Where exploration starts & Allocation and coordination & What determines subsequent expansion \\
\midrule
C & Retained main-search state & All remaining budget to one route & Existing frontier and adjacency cursors \\
MEP-$m$ & Up to $m$ backup upper entrances & Paid descent endpoints enter a shared queue; $m=1,4,8,16$ & Common queue order; no private scout interval \\
FRONTIER & Up to four pending main-frontier candidates & F's local scouts, with cached roots and shared progress & Unfinished scout work returns to the main frontier \\
FANNG search & Common upper-descent endpoint & One edge at a time under the distance budget & Query-prioritized source and remaining-edge cursor \\
iQAN~\cite{iqan} & Multiple candidates from the search pool & Local queues, staged path expansion, adaptive synchronization & Local paths feed candidates back into global search \\
CDSM-F & Backup entrances after a main-search prefix & Up to four local scouts; 400 new distances each & Active, deferred, and current scout work joins the main frontier with shared cursors \\
\bottomrule
\end{tabularx}
\end{table*}

